\documentclass[10pt]{article}

\usepackage[utf8]{inputenc}
\usepackage[T1]{fontenc}
\usepackage[spanish]{babel}

\usepackage{aleph-comandos}
\usepackage{moodle}
% Se recomienda leer la documentación de estos
% paquetes en https://www.alephsub0.org/recursos/

% -- Paquetes adicionales
\usepackage{amssymb}
\usepackage{mathpazo}

% -- Datos
\title{Reactivos Álgebra Lineal}
\author{Evelyn Cueva}

\usepackage[margin=2cm]{geometry}

\begin{document}
\maketitle

\begin{quiz}{Algebra-Lineal-Examen-Media-Carrera-2023B}

% Ejercicio de cálculo para evaluar núcleo e imagen de una transformación lineal. Matriz transpuesta y composición de transformaciones 

\begin{multi}{algebra-lineal-2023B-01}
    Considere $T$ una tranformaci\'on lineal que transpone cualquier matriz de $2\times 2.$
    Indique cu\'al de las siguientes afirmaciones es falsa:
    \item[fraction=0] $T^2$ es la transformaci\'on identidad. 
    \item[fraction=0] El n\'ucleo de $T$ es la matriz cero. 
    \item[fraction=0] Toda matriz de $2\times 2$ est\'a en la imagen de $T.$
    \item* Si $M$ es una matriz de $2\times 2$, $T(M)=-M$ no puede ocurrir. 
\end{multi}

% Pregunta conceptual para evaluar las propiedades basicas de determinantes. 

\begin{multi}{algebra-lineal-2023B-02}
    Sean $A,B,C$ matrices de tama\~no $n\times n$. Si $|A|$ denota el determinante de $A$, identifique cu\'al de las siguientes afirmaciones es verdadera: 
    \item[fraction=0] El determinante de $I+A$ es $1+|A|.$
    \item[fraction=0] Sean $A,B,C$ matrices de $n\times n$. El determinante de $ABC$ es $|A||B||C|,$ 
    \item[fraction=0] El determinante de $4A$ es $4|A|.$
    \item* El determinante de $AB-BA$ siempre es cero.  
\end{multi}

% Pregunta conceptual para evaluar conocimientos de tranformaciones lineales respecto de las dimensiones del espacio vectorial, espacio nulo e imagen. 

\begin{multi}{algebra-lineal-2023B-03}
Sean $V$ y $W$ dos espacios de dimensi\'on finita, $\mathcal L(V,W)$ el espacio de todas las transformaciones lineales de $V$ en $W$. Para $T\in\mathcal L(V,W)$, notamos por $\mathcal N(T)$. Seleccione la afirmaci\'on que no es siempre verdadera:
     \item[fraction=0]  $T$ es inyectiva si y solo si $\mathcal N(T)=\{0\}.$
     \item* Si $\dim V$ es menor que $\dim W,$ entonces $T$ no es inyectiva. 
     \item[fraction=0] Si $\dim V$ es menor que $\dim W,$ entonces $T$ no es sobreyectiva. 
     \item[fraction=0] $T(0)=0.$
\end{multi}

% Pregunta conceptual y de calculo para evaluar independencia lineal y espacios generados por listas de diferentes tamaños.

\begin{multi}{algebra-lineal-2023B-04}
De los siguientes conjuntos de vectores en $\Rbb^3$, determine cu\'al no genera todo el espacio $\R^3$.
    \item[fraction=0]  $\left\lbrace(1,-1,1),(-1,-1,1),(0,1,1)\right\rbrace$
    \item[fraction=0]  $\left\lbrace (1,-1,1),(-1,-1,1),(0,1,1),(0,-1,3) \right\rbrace$
    \item*  $\left\lbrace(8,-3,4),(0,3,0),(8,0,4)\right\rbrace$
    \item[fraction=0] $\left\lbrace (1,0,0),(0,1,0),(0,0,1) \right\rbrace$
\end{multi}

% Pregunta conceptual para identificar la relacion entre propiedades de matriz inversa, sistemas de ecuaciones y sistemas homogeneos

\begin{multi}{algebra-lineal-2023B-05}
Sea $A$ una matriz de tama\~no $n\times n$. De las siguientes afirmaciones, seleccione aquella que no sea equivalente con las dem\'as:
    \item[fraction=0] $A$ es no singular.
    \item[fraction=0] El sistema $Ax=0$ tiene soluci\'on \'unica trivial.
    \item* $A$ no es equivalente por filas con la matriz identidad. 
    \item[fraction=0] El determinante de $A$ es no nulo.
\end{multi} 
\end{quiz}

\end{document}